\section{Incident Command：一条外部报告何时升级为企业事件}

当报告显示 credential 被使用、敏感系统被访问、数据可能离开控制边界或攻击仍在持续，组织面对的已经不只是 vulnerability remediation。\term{incident commander}（事件指挥官，IC）是负责协调目标、优先级、节奏、角色与升级的单一运行负责人；IC 不必亲自完成取证或作出法律结论，但要确保每个 workstream 有 owner，冲突被升级，决策被记录。

\subsection{跨职能指挥结构}

事件响应至少包含 command、technical、legal/comms 和 business 四个 workstream。Technical 团队确认事实、控制影响、修复并恢复；legal/comms 识别义务、privilege 和外部受众；business 处理用户、运营、财务和产品取舍；command 维护共同 timeline、objectives、cadence 与 dependency。让 security on-call 同时担任所有角色，会在压力下把假设误写成事实。

\lecturefigure{06-incident-command.png}{重大安全事件需要 command、technical、legal/comms 与 business 的跨职能结构。}{NIST SP 800-61 Rev. 3；本地字幕 15:10--25:40；概念重绘。}

\begin{knowledgebox}{读图：IC 负责同步，不垄断专业判断}
Command 维护 severity、owners 和 decision log；technical 处理 contain、forensics、recovery；legal/comms 分析 privilege、regulators、customers 与 messaging；business 决定 customer impact、continuity 和资源。箭头的含义是双向：技术发现会改变 materiality，法律 deadline 会改变调查优先级，业务恢复动作可能破坏证据，因此任何工作流都不能独立运行。
\end{knowledgebox}

\begin{knowledgebox}{术语消化：severity、containment 与 recovery}
\term{Severity} 是组织对当前影响和紧迫度的运行分类，不等同最终法律 materiality。\term{Containment} 是限制攻击或数据影响的短期动作，例如隔离 credential 或服务。\term{Eradication} 消除根因与 persistence；\term{recovery} 恢复可信服务并监控复发。先恢复业务并不意味着调查结束，先保全证据也不意味着必须延迟所有保护用户的动作。
\end{knowledgebox}

\teachervoice{课堂叙述中最值得保留的不是具体入侵步骤，而是分类冲突：一条线索可以最初由 disclosure channel 接收，却逐步显示为更严重事件。实践经验是预先定义升级条件，让员工不必靠“这看起来像不像 bounty”作个人判断。}

\begin{importantbox}{VDP-to-incident escalation triggers}
发现真实账户或 credential 使用、非最小化数据访问、敏感数据离开边界、持续访问、第三方受影响、研究者要求与停止损害无关的付款、事实与报告明显不一致，或任何监管/合同 trigger 时，立即启动 incident review。触发不等于指控研究者犯罪，而是提高组织调查和治理等级。
\end{importantbox}

\subsection{本章小结}

事件指挥把多个专业判断组织成共同运行节奏。IC 的价值是让事实、保护动作、法律义务和经营决定互相可见，而不是由 CSO 一人扮演工程师、律师、发言人和 CEO。

